\ConclusionGenerale

La conclusion générale doit répondre aux questions posées dans l’introduction. Elle permet de faire un bilan des résultats obtenus, d’indiquer les limites du travail et d’ouvrir sur des perspectives futures.

Structure recommandée :

\textbf{Rappel du contexte et des objectifs}
\begin{itemize}
	\item Résumé rapide du problème traité et des objectifs fixés.
\end{itemize}

\textbf{Synthèse des résultats}

\begin{itemize}
	\item Présentation des principales contributions du mémoire.
	\item Validation (ou non) des hypothèses initiales.
\end{itemize}

\textbf{Limites et difficultés rencontrées}

\begin{itemize}
	\item Discussion des éventuelles contraintes ou limites du travail.
\end{itemize}

\textbf{Perspectives et recommandations}

\begin{itemize}
	\item Proposition d’améliorations ou de pistes pour des travaux futurs. 
\end{itemize}